\chapter{L'oscilloscope}

\section{Une voie par voix}
L'oscilloscope dessine la forme du son de chaque voix. Chaque puce occupe une colonne, chaque voix une ligne, avec
son nom en haut de la case. Vous voyez ainsi qui joue quoi : la basse qui tient, l'arpège qui scintille, la caisse
claire qui passe en un éclair.

Pour que les formes restent immobiles, nous calons la fenêtre de chaque voie sur la période de sa note : c'est le
déclenchement, comme sur un oscilloscope de laboratoire. Nous réduisons ensuite le signal à deux points par colonne
de pixels, le minimum et le maximum, si bien qu'aucun pic n'est perdu même quand la voie est étroite.

\section{Cas particuliers}
Avec la voie des samples du SID activée, une quatrième ligne par puce montre les échantillons joués par le volume.
Pour l'AdLib, nous disposons les neuf voix mélodiques en 3 sur 3 et masquons les cinq voies du mode rythme, que peu
de morceaux utilisent ; activez « AdLib rhythm lanes » si vous en avez besoin.

\section{La réponse du filtre du SID}
Le bouton de visualisation remplace la moitié haute de l'oscilloscope par la courbe du filtre de chaque SID, avec la
position de chaque voix. Les états précédents restent visibles quelques instants en fondu, ce qui rend lisibles les
balayages rapides de la coupure.
\capturedouble{ux4-lowpass-waves}{ux4-sweep}{Filtre passe-bas résonant}{balayage de la coupure}
Si vous voulez plus de place pour les voix, réduisez la part de la réponse du filtre dans les réglages de
l'oscilloscope.
\capture[0.7]{ux4-lowpass-share35}{Réponse du filtre réduite à 35~\% de la hauteur}

\section{Réglages}
\capture{prefs-en-oscilloscope}{Réglages de l'oscilloscope}
Vous réglez la fenêtre de temps, le gain, le nombre de périodes affichées, la hauteur de l'oscilloscope, la part de la
réponse du filtre, la fréquence d'images maximale et le moteur de rendu. Vulkan donne le rendu le plus fluide quand
votre téléphone le permet ; sinon, OpenGL ES prend le relais.
